\section{Uncertainty quantification: constructing the covariance blocks}
\label{sec:systematics}



The uncertainty budget combines three independent blocks --- systematic
(universe), statistical (bootstrap), and ML (training-seed variation) --- summed assuming
independence (distinct RNGs / physics sources). The 3D and 4D budgets
(\S\ref{sec:3d-syst}, \S\ref{sec:4d}) repeat the identical procedure on the
higher-dimensional grids. Full mathematical procedures, with worked examples
on the reported numbers, are in App.~\ref{app:statmethods}. This section describes
how each block is built; the resulting budget, and its agreement with the published
total, are reported with the measurement in \S\ref{sec:results-uq}.


\subsection{Systematic universes}
\label{sec:uq-universes}
The detector and model systematics use the MAT-MINERvA universe machinery: 44
error bands and 187 universes in total, identical to the published analysis's
\texttt{GetStandardSystematics} set. Each universe is re-unfolded with the
central estimator seed. All per-band covariances use the MAT convention:
universe-mean centering and biased $1/N$. For a $\pm1\sigma$ pair this is the
outer product of the half-difference; a common displacement of both endpoints
from CV is excluded from the variance. No per-band displacement artifact is
stored; only the joint-throw mean shift described below is stored separately. The 100 PPFX
universes use the same mean-centered $1/N$ formula.

Dump-all shifted branches are limited to CV event-loop support. For the five
genuinely kinematic bands (BeamAngleX/Y, MuonResolution, and
Muon\_Energy\_MINERvA/MINOS), selection-complete active-universe event loops
are the production migration gate and are still in progress; until their
replacement covariance lands, higher-dimensional uncertainty products that use
the dump-all lateral branches remain candidates. MinosEfficiency and GEANT are
weight-only. Top 2D bands: Flux \SI{\fluxBand}{\percent}, Muon\_Energy\_MINOS
\SI{2.31}{\percent}, MinosEfficiency \SI{1.47}{\percent},
Muon\_Energy\_MINERvA \SI{1.29}{\percent}. Figure~\ref{fig:uqbands} shows the
grouped fractional-uncertainty projections in the published Fig.~6/7 style
(the full $205\times205$ covariance projected exactly, $C_{1D}=PC_{2D}P^{\top}$,
before dividing by the 1D central value).

\begin{figure}[H]
  \centering
  \includegraphics[width=0.49\textwidth]{MEFHC_fig6_7_uncertainty_pt}
  \includegraphics[width=0.49\textwidth]{MEFHC_fig6_7_uncertainty_pz}
  \caption{Fractional uncertainty on the 2D cross section projected onto $\pT$
  (left) and $\pPar$ (right), grouped by source in the style of Figs.~6--7 of
  Ref.~\cite{MINERvA:2021owq}: flux, interaction models, normalization,
  hadronic response, muon reconstruction, and the statistical band.}
  \label{fig:uqbands}
\end{figure}

\paragraph{Flux $1/\Phi$ normalization.}
The dominant flux term enters as $\sigma\propto1/\Phi$; varying $\Phi$ per flux
universe, rather than dividing every universe by the CV $\Phi$, gives a Flux
band of \SI{\fluxBand}{\percent} (versus $\sim\SI{1}{\percent}$ under the CV
convention) and a standalone combined budget of \SI{\uqMedian}{\percent}. The \SI{1.4}{\percent} target-nucleon
normalization is a fully-correlated rank-1 band. The 100-universe flux
covariance has a participation-ratio effective rank of $\approx1$ (one coherent
normalization mode carries \SI{99.6}{\percent} of the variance), and its leading
mode correlates at $\rho=\rhoFluxMuon$ with the rank-1 muon-energy-scale band:
the rederived Bashyal~\cite{Bashyal:2021tqu} joint block, whose cross term is
rank-2, with two independent estimators giving $\rho=0.87$--$0.90$. Adding the
block leaves the systematic rank unchanged --- its directions already lie within
the span of the universe covariance (App.~\ref{sec:rank}).

\subsection{Statistical uncertainty: Poisson bootstrap}
\label{sec:uq-boot}
The statistical block is a 300-replica per-event Poisson(1) weight bootstrap
applied jointly to data and simulation, with two independent sub-RNGs (data
and MC are independent variance sources), the MC reco and truth weights riding
a \emph{single} per-event draw (preserving the within-event migration
correlation), and the GBDT seed pinned per replica so the bootstrap and
training-seed variances are separable. Each replica is a full re-unfold; the
covariance is the sample covariance over replicas
(Fig.~\ref{fig:bootstrap}). Our bootstrap block is genuinely smaller than the
published statistical block ($\approx2.5\times$ by trace). We report this as an
\emph{observed} difference in the statistical covariance rather than a
demonstrated causal efficiency advantage; it is consistent with an unbinned
statistical-efficiency gain, and the data-only stream of the data/MC split
bootstrap ($C_{\mathrm{data}}/\texttt{StatOnly}=0.356$) confirms it is not a
missing term (App.~\ref{sec:rank}).

\begin{figure}[H]
  \centering
  \includegraphics[width=0.49\textwidth]{uq_spread_2d}
  \includegraphics[width=0.49\textwidth]{uq_corr_2d}
  \caption{300-replica Poisson bootstrap: per-bin relative spread of the
  unfolded 2D cross section (left) and the bootstrap correlation matrix on the
  205 reported bins (right).}
  \label{fig:bootstrap}
\end{figure}

\subsection{ML/optimization uncertainty: training-seed variation}
\label{sec:uq-ml}
ML stochasticity is budgeted by re-unfolding with 10 independent estimator
seeds at fixed data and fixed universe (Fig.~\ref{fig:seedscan}); the sample
covariance over seeds is the ML block. It is the smallest block
(median \SI{0.166}{\percent} per bin). Two refinements harden it: (i) a
\emph{train/test-split} variant, which also re-draws the internal training
split, gives a band $1.24\times$ the pure-seed band, moving the combined 3D
budget by only $+0.04\%$ in $\sqrt{\mathrm{Tr}\,C}$ --- the published budgets use
the split-inclusive band; (ii) the ensemble-mean central value (averaging the
seed ensemble) agrees with the production single-run CV to \SI{0.28}{\percent}.
For the corrected 5D candidate campaign, the ML ensemble fixes the estimator seed
at 42 and varies the train/test split seed. It therefore measures split response,
not the full estimator stochasticity. A dedicated fixed-data estimator-seed scan
(twelve estimator seeds at fixed data and fixed split) closes this scope: its
covariance matrix has $\sqrt{\mathrm{Tr}\,C}=\SI{\gbdtAiEstTrace}{cm^2/nucleon}$,
\SI{\gbdtAiEstFrac}{\percent} of the candidate split-varied band
(\SI{\gbdtMlSplitTrace}{cm^2/nucleon}), so the adopted band is dominated by
genuine estimator stochasticity rather than by the particular train/test
partition. The scan remains an auxiliary robustness check rather than an
independent matrix --- split and estimator responses are not assumed
uncorrelated --- and both bands are minor components, about three percent of
the combined budget in $\sqrt{\mathrm{Tr}\,C}$.

\begin{figure}[H]
  \centering
  \includegraphics[width=0.49\textwidth]{seedscan_spread_2d}
  \includegraphics[width=0.49\textwidth]{seedscan_band_pt}
  \caption{ML training-seed variation: per-bin relative spread over 10 estimator seeds (left)
  and the resulting fractional band projected onto $\pT$ (right).}
  \label{fig:seedscan}
\end{figure}

\subsection{Combination and the unified-throw test}
\label{sec:uq-combine}
The nominal covariance decomposition is the block sum
$C = C_{\mathrm{syst}} + C_{\mathrm{stat}} + C_{\mathrm{ML}}$, which assumes
the blocks are independent \emph{and} that the unfolding responds linearly to
simultaneous perturbations (so that per-band covariances add). The second
assumption is tested directly by the \emph{unified throw}
(\S\ref{sec:4d}): joint re-unfolds with all reweightable knob bands and the
flux multisim perturbed simultaneously per event.  The corrected construction
uses the actual $-1\sigma$ and $+1\sigma$ event-weight endpoints, interpolates
asymmetrically between them, holds the unfolding-estimator seed fixed for every
throw, and forms the joint covariance about the throw mean.  The joint
displacement of that mean from the CV is stored separately as a per-bin vector
and its norm; per-band endpoint displacements are not stored. The matched block comparator
uses the MAT convention: universe-mean centering and the biased $1/N$
normalization, with both endpoints re-unfolded.

For the 5D candidate this construction was carried out both mean-centered and
CV-centered, with the joint mean shift reported separately rather than folded into
the covariance. \textbf{No magnitude from either construction is quoted here.} Both are
pre-J28 and background-aware, both sit under the 2026-07-12 construction-cause
quarantine, and the corrected J28 products are footed on the
\emph{non}-background-aware configuration --- so no substitution is available, and
supplying one would make the sentence false rather than current. Neither is adopted for
publication until the selection-complete lateral replacement lands. No
positive-semidefiniteness measurement exists for either: \texttt{VL16} and \texttt{VL17}
record that property for the \emph{successor} background-aware products only, and it does
not transfer by inheritance. Routed record:
\texttt{docs/orchestration/PROCEDURE-gbdtFive-macro-update.md}.

Repeating all 188 universe unfolds with
per-universe background subtraction changes the combined \emph{block-sum}
sqrt-trace by only \SI{0.28}{\percent}; the effect on the \emph{adopted}
covariance is smaller still, \SI{0.09}{\percent} before the flux (J28)
correction and \SI{0.18}{\percent} after it, because the adoption's per-bin
$\max$ inflation transfer damps the block-sum change. These are ratios taken against
the superseded products described above, so they bound the size of the background-CV
choice rather than describing a current covariance; either way that concern is
numerically negligible. The 4D and extended-fiducial results do not yet have covariances
constructed with this procedure; the 5D covariance is not transferred to them.

More generally, this reframes the block sum as an approximation --- valid only
where the unfolding responds linearly --- rather than as a default
construction.  We retain the per-source block decomposition for attribution,
but use it only after joint throws and matched block endpoints pass the stated
validation tests.

\paragraph{$\chi^{2}$ caveat.}
The combined paper$+$ours $\chi^{2}/\mathrm{ndf}=1.481$ double-counts systematics
shared with the paper (flux/GENIE/MINOS); the ours-only inverse-covariance
$\chi^{2}$ is evaluated only with a truncated-spectral convention on the
well-constrained subspace. We therefore quote per-bin pulls, the regularization
scan, and the truncated result together rather than treating an unconstrained
pseudo-inverse as a standalone goodness-of-fit (App.~\ref{sec:rank}).
